\scanpage{226}
% slanted
\margin{cette formule a un sens et est \unsure{vraie|vrai} sans supposer $\mathcal B$ karoubienne}

\[
 F(x)\simeq\operatorname{Hom}_k^*(\alpha^*(x),\underset{\substack{\rotatebox{90}{$=$}\,\mathrm{déf}\\ F(\mathbb C^*)}}{C_F^*})\quad\deleted{\unsure{*k}}
\]
où la définition du $\operatorname{Hom}^*$ se fait comme dans la prop. 3.

\deleted{\diagram{\begin{tikzcd}[ampersand replacement=\&,column sep=small]
 F(x)\otimes\alpha^{[n]}(x)\arrow[r]\& C_F^{[n]}\arrow[d,equal]\\
 \operatorname{Hom}(x,C^{[n]})\arrow[u,equal]\& F(\mathbb C^{[n]})
\end{tikzcd}}
% slanted
\margin{\unsure{par définition|par définitions} de $\alpha^*$ \deleted{et $DP$} $\alpha^*$ est \deleted{\unsure{de $DP$}} la commutation de $DP$ aux foncteurs \unsure{add.|std.|etc.}}
On procède comme pour la prop. 3 pour définir la flèche dans $\widehat{\mathcal B}$, et pour montrer que c'est un iso, on se ramène au cas $x$ de la forme $\mathbb C^{i}_{\unsure{*|\bullet|\xi}}$, il faut prouver qu'alors
\[
 F(\mathbb C^i)=C^i\longrightarrow\operatorname{Hom}_k^*(\unsure{\alpha^*},\mathbb C^\bullet)
\]}

En effet, on trouve \deleted{\unsure{tout}} de suite que si $K^\bullet\in\mathrm{Ab}_k^\bullet$, $C^\bullet\in\mathcal B^\bullet$, $K^*=DP(K^\bullet)$ et $C^*=DP(K^\bullet)$, alors on a un iso. can. bifonctoriel dans $\widehat{\mathcal B}$
\[
 DP\colon\operatorname{Hom}_k^\bullet(K^\bullet,C^\bullet)\simeq\operatorname{Hom}_k^*(K^*,C^*),
\]
et on conclut par prop 3.

\deleted{\uline{Prop 6} \textcircled{*} Le foncteur
\[
 ND\colon\mathrm{Ab}_{k*}\longrightarrow\mathrm{Ab}_{k\bullet}
\]
est donné par
\[
 (ND)(L_*)\simeq\operatorname{Hom}_{k,*}(\unsure{\mathbb C|C}_*^\bullet,L_*)
\]
i.e.
\[
 ND(L_*)_i\simeq\operatorname{Hom}_{k,*}(\underset{\substack{\rotatebox{90}{$=$}\\\Lambda^{i\unsure{+1}}\Phi_*}}{C_*^i},L_*)\quad\text{pour }i\geq0
\]
l'hom \inserted{bord} $ND(L_*)_i\to ND(L_*)_{i-1}$ $(i\geq1)$ étant le transposé de $C_*^{i-1}\longrightarrow C_*^i$.}

% archive sheet number: 225

